\documentclass[twocolumn]{iopjournal}

\usepackage{amsmath,amssymb,amsfonts}
\usepackage{array}

\fancyhfoffset[L]{0mm}
\setlength{\marginparwidth}{0pt}
\setlength{\marginparsep}{0pt}

% ==========================================
% ZENODO PREPRINT HEADERS (ACTIVE)
% ==========================================
\renewcommand{\articletype}[1]{%
  {\vspace*{-8mm}\noindent \Large \sf Preprint}\par
  \vspace*{6mm}\noindent{\scriptsize \sf{\bfseries \MakeUppercase{#1}}}}

\fancyhead[L]{}
\fancyhead[R]{}

\newcommand{\sgn}{\operatorname{sgn}}

\hypersetup{pdftitle={The Two Roots of the Tetrad Equation}, pdfauthor={Stephen Nelson}, pdfsubject={Gravity Field Dynamics (GFD)}, pdfkeywords={tetrad, scheme, Lorentz group, root-blindness, spatial parity, stellated octahedron, proof}}

\begin{document}

\articletype{Supplementary Material}

\title{The Two Roots of the Tetrad Equation}

\author{Stephen Nelson}

\affil{Field Dynamics Research}

\ead{stephen@fielddynamics.org}

\begin{abstract}
The tetrad equation $g = e^T \eta\, e$ forces
$(\det e)^2 = \det g / \det \eta$, an equation of the form
$y^2 = c$. We prove that its scheme has exactly two connected
components, one over each root $\det e = \pm\sqrt{c}$, and that
neither splits further over $\mathbb{C}$. A change of frame exchanges
the roots if and only if it reverses orientation; in three spatial
dimensions spatial parity does, and it is the only such map that
commutes with every rotation and reflection of an observer's space.
For future-pointing frames the two roots are the two tetrahedra of the
stellated octahedron. The metric determines $(\det e)^2$ and never
$\det e$, so keeping one root, as the restriction to proper
orthochronous frames does, is a convention.
\end{abstract}

\keywords{tetrad, scheme, Lorentz group, root-blindness, spatial parity, stellated octahedron, proof}

% ----------------------------------------------------------------------

\section*{Setting}

At a point of a spacetime with $1 + D$ dimensions and metric $g$ of
signature $(-, +, \dots, +)$, a \emph{tetrad} is a real square matrix
$e = (e^a{}_\mu)$ with
\begin{equation}
  e^T \eta\, e = g, \qquad \eta = \mathrm{diag}(-1, I_D).
  \label{eq:tetrad}
\end{equation}
The entries of $e$ are the unknowns, and Eq.~\eqref{eq:tetrad} is a
set of quadratic equations in them. They define the affine scheme
$X_g = \operatorname{Spec} \mathbb{R}[e^a{}_\mu]/(e^T \eta\, e - g)$
over $\mathbb{R}$~\cite{Grothendieck1960,Hartshorne1977}, whose real points $X_g(\mathbb{R})$ are the real tetrads. The
orthogonal group scheme $O(\eta)$, with real points the Lorentz group
$O(1,D)$, acts by $e \mapsto \Lambda e$. We write
$P = \mathrm{diag}(1, -I_D)$ (spatial parity),
$T = \mathrm{diag}(-1, I_D)$ (time reversal) and
$R = \mathrm{diag}(1, -1, 1, \dots, 1)$ (one spatial reflection).
Since $g$ and $\eta$ have the same signature,
$c = \det g / \det \eta > 0$.

% ----------------------------------------------------------------------

\section*{Two Lemmas}

\paragraph{Lemma 1 (time cone).} \textit{If $\Lambda \in O(1,D)$ and
$\eta(u,u) < 0$, then
$\sgn (\Lambda u)^0 = \sgn \Lambda^0{}_0 \cdot \sgn u^0$, and
$\Lambda \mapsto \sgn \Lambda^0{}_0$ is a homomorphism.}

\paragraph{Proof.} $\Lambda^T \eta \Lambda = \eta$ gives
$\Lambda \eta \Lambda^T = \eta$, whose $(0,0)$ entry is
$|r|^2 = (\Lambda^0{}_0)^2 - 1$ for $r = (\Lambda^0{}_i)$. A timelike
$u$ has $|\vec u| < |u^0|$, so in
$(\Lambda u)^0 = \Lambda^0{}_0 u^0 + r \cdot \vec u$ the second term
is smaller in modulus than the first. Taking $u = \Lambda'_{\cdot 0}$
gives the homomorphism. \hfill $\square$

\paragraph{Lemma 2 (centre).} \textit{For $D \ge 1$, the elements of
$O(1,D)$ commuting with all of $O(1,D)$ are $\pm I$, and those of
$O(D)$ commuting with all of $O(D)$ are $\pm I_D$.}

\paragraph{Proof.} Write $Z = (\alpha, w^T; v, Y)$. Commuting with
$\mathrm{diag}(1, -I_D)$ gives $v = w = 0$; with all diagonal sign
matrices and all permutation matrices in $O(D)$, $Y = \lambda I_D$,
which also settles $O(D)$. A boost in the $0$--$1$ plane forces
$\alpha = \lambda$, and $Z^T \eta Z = \lambda^2 \eta$ gives
$\lambda = \pm 1$. \hfill $\square$

% ----------------------------------------------------------------------

\section*{The Two Roots}

\paragraph{Theorem 1 (the scheme).} \textit{Let $e_0$ be a real
tetrad. Then:}
\begin{itemize}
\item[(a)] \textit{$\Lambda \mapsto \Lambda e_0$ is an isomorphism
  $O(\eta) \to X_g$. Every real tetrad is $\Lambda e_0$ for exactly
  one $\Lambda \in O(1,D)$.}
\item[(b)] \textit{On $X_g$, $(\det e)^2 = c$. So
  $\det : X_g \to \{y^2 = c\}$, and $X_g = X_+ \sqcup X_-$, where
  $X_\pm$ is the fibre over $y = \pm\sqrt c$. Each $X_\pm$ is
  connected, and stays connected over $\mathbb{C}$: $X_g$ has exactly
  two connected components, the two roots.}
\item[(c)] \textit{$e \mapsto \Lambda e$ exchanges the roots if and
  only if $\det \Lambda = -1$. $T$ and $R$ always do; $P$ does if and
  only if $D$ is odd; $-I$ does if and only if $D$ is even.}
\end{itemize}

\paragraph{Proof.} (a) From Eq.~\eqref{eq:tetrad},
$(\det e)^2 \det \eta = \det g$, a unit, so $e$ is invertible on
$X_g$ and $e \mapsto e\,e_0^{-1}$ is the inverse map: it lands in
$O(\eta)$ because $e_0^{-T} e^T \eta\, e\, e_0^{-1} = e_0^{-T} g\,
e_0^{-1} = \eta$. (b) The identity above divided by $\det \eta$. The
target is two points, so $X_\pm$ are open and closed. Under (a),
$X_+ \cup X_-$ is $O(\eta) = SO(\eta) \sqcup R\,SO(\eta)$, and
$SO(\eta)$ is connected over $\mathbb{C}$, where it is
$SO(1 + D, \mathbb{C})$~\cite{Borel1991}. (c)
$\det(\Lambda e) = \det \Lambda \det e$, with $\det T = \det R = -1$,
$\det P = (-1)^D$ and $\det(-I) = (-1)^{D+1}$. \hfill $\square$

\paragraph{Theorem 2 (real points and parity).} \textit{Fix a
future-directed timelike $t$ and set $\tau(e) = \sgn (e\,t)^0$.}
\begin{itemize}
\item[(a)] \textit{$\tau(\Lambda e) = \sgn \Lambda^0{}_0 \cdot
  \tau(e)$. The real points of each root form two pieces,
  $\tau = \pm 1$; $X_g(\mathbb{R})$ has four.}
\item[(b)] \textit{Among future-pointing frames, $\tau = +1$, the
  root $\sgn \det e$ is the orientation of the spatial legs. A change
  of frame that keeps the observer, $\Lambda = \mathrm{diag}(1, Q)$,
  exchanges the roots if and only if $\det Q = -1$, and for $D$ odd
  the only such $Q$ that commutes with every rotation and reflection
  of the observer's space is $-I_D$: spatial parity $P$.}
\item[(c)] \textit{Other than the identity, $e \mapsto -e$ is the
  only change of frame that commutes with all others. For $D$ odd it keeps the root and reverses
  $\tau$: $-e = T(Pe)$.}
\end{itemize}

\paragraph{Proof.} (a) $u = e\,t$ has $\eta(u,u) = g(t,t) < 0$; apply
Lemma~1. Under Theorem~1(a), $\tau$ and the root separate the four
components $X \cdot SO^+(1,D)$, $X \in \{I, T, R, RT\}$, of
$O(1,D)$ (for $D = 3$, see~\cite{Weinberg1972}). (b) $\det \mathrm{diag}(1, Q) = \det Q$;
the centre of $O(D)$ is $\pm I_D$ (Lemma~2), and
$\det(-I_D) = -1$ for $D$ odd. (c) Lemma~2, with
$\det(-I) = (-1)^{D+1}$ and $(-I)^0{}_0 = -1$. \hfill $\square$

\paragraph{Theorem 3 (root-blindness).} \textit{The metric determines
$(\det e)^2 = c$ and never $\det e$. Every quantity built from $g$,
from the connection and curvature to the Einstein--Hilbert action,
takes one value on both roots, for
tetrad fields as at a point. Restricting frames to
$SO^+(1,D)\,e_0$, the usual choice for oriented and time-oriented
frames~\cite{Wald1984,Hehl1976}, keeps one root; $g$ does not
determine which.}

\paragraph{Proof.} By Theorem~1(a) every point of $X_g$ gives the same
$g$, and $Pe(x)$ solves Eq.~\eqref{eq:tetrad} wherever $e(x)$ does.
$SO^+(1,D) \subset SO(\eta)$, and $e_0$ may be taken in either root.
\hfill $\square$

% ----------------------------------------------------------------------

\section*{The Two Tetrahedra}

Take $D = 3$. Let $E_a$ be the frame vectors of $e$, with
$g(E_a, E_b) = \eta_{ab}$, and $W$ the rest space spanned by
$E_1, E_2, E_3$. Letter~1's directions~\cite{Nelson2026a} are
$n_a = \sum_i M_{ia} E_i$, with
\begin{equation}
  M = \begin{pmatrix}
    +1 & +1 & -1 & -1 \\
    +1 & -1 & +1 & -1 \\
    +1 & -1 & -1 & +1
  \end{pmatrix},
  \label{eq:M}
\end{equation}
the regular tetrahedron forced by the tetrad~\cite{NelsonForced}.

\paragraph{Theorem 4.} \textit{In the basis $E_i$ of $W$:}
\begin{itemize}
\item[(a)] \textit{$n(Pe) = n(-e) = -M$ and $n(Te) = M$. A
  future-pointing frame of the same observer, $\mathrm{diag}(1, Q)\,e$,
  gives $n = Q^T M$: a rotation of $M$ in the root of $e$, a rotation
  of $-M$ in the other.}
\item[(b)] \textit{No $Q \in SO(3)$ has $QM = -M$. For a regular
  simplex $M_D$ in $\mathbb{R}^D$, such a rotation exists if and only
  if $D$ is even.}
\item[(c)] \textit{The region
  $\mathrm{TetA} = \{x \in W : \max_a n_a \cdot x \le 1\}$ of $e$
  becomes $-\mathrm{TetA} = \{\max_a (-n_a \cdot x) \le 1\}$ for $Pe$,
  the all-flipped region TetB of Letter~1; together they form the
  stellated octahedron~\cite{NelsonForced}.}
\end{itemize}

\paragraph{Proof.} (a) The frame vectors of $\Lambda e$ are
$E\Lambda^{-1}$; for $\mathrm{diag}(1, Q)$ the spatial ones become
$E_j' = \sum_i Q_{ji} E_i$, and $Q^T M = -(-Q^T)M$ with
$\det(-Q^T) = -\det Q$. (b) $\mathrm{rank}\, M_D = D$, so
$QM_D = -M_D$ forces $Q = -I_D$, of determinant $(-1)^D$. (c)
$(-n_a) \cdot x = n_a \cdot (-x)$. \hfill $\square$

\paragraph{Remark.} The tetrahedron is congruent to its mirror image:
a $90^\circ$ rotation maps the vertex set of TetA onto that of TetB,
but reorders the legs. The roots differ in the order of the legs and
the region of $W$ they occupy, not in shape. And the tetrahedron
follows the root only for future-pointing frames: $-e$ carries TetB
in the root of $e$, because it also reverses time.

% ----------------------------------------------------------------------

\section*{Scope}

\vspace{2pt}
\noindent
\begin{tabular*}{\columnwidth}{@{\extracolsep{\fill}}l>{\raggedright\arraybackslash}p{0.66\columnwidth}@{}}
\hline
Ingredient & Without it \\
\hline
$\det g \ne 0$ & $c = 0$: a double root, not two \\
$D$ odd & $P$ keeps the root; for $D = 2$, $-M_2$ is $M_2$ turned by
  $\pi$ \\
Future-pointing & $-e$ carries TetB in the same root \\
Ordered legs & As point sets TetB is a rotation of TetA \\
\hline
\end{tabular*}
\vspace{4pt}

% ----------------------------------------------------------------------

\begin{thebibliography}{9}

\bibitem{Grothendieck1960}
A.~Grothendieck and J.~Dieudonn\'{e},
\textit{\'{E}l\'{e}ments de g\'{e}om\'{e}trie alg\'{e}brique~I},
Publ.\ Math.\ IH\'{E}S \textbf{4} (1960).

\bibitem{Hartshorne1977}
R.~Hartshorne,
\textit{Algebraic Geometry}
(Springer, New York, 1977).

\bibitem{Borel1991}
A.~Borel,
\textit{Linear Algebraic Groups}, 2nd ed.\
(Springer, New York, 1991).

\bibitem{Weinberg1972}
S.~Weinberg,
\textit{Gravitation and Cosmology}
(Wiley, New York, 1972).

\bibitem{Wald1984}
R.~M.~Wald,
\textit{General Relativity}
(University of Chicago Press, Chicago, 1984).

\bibitem{Hehl1976}
F.~W.~Hehl, P.~von der Heyde, G.~D.~Kerlick and J.~M.~Nester,
Rev.\ Mod.\ Phys.\ \textbf{48}, 393 (1976).

\bibitem{Nelson2026a}
S.~Nelson,
\textit{The Coordinate-Space Geometry of the Tetrad Frame},
Field Dynamics Research (2026).

\bibitem{NelsonForced}
S.~Nelson,
\textit{The Spatial Geometry of the Lorentz Tetrad},
Field Dynamics Research (2026).

\end{thebibliography}

\end{document}
